\chapter{Photogrammetry, Depth Fusion, and Apple Capture APIs}
\label{ch:theory1}

This chapter explains the geometric and algorithmic ideas behind the capture pipelines used in the thesis application. It first reviews photogrammetry and the classical chain structure from motion (SfM) $\rightarrow$ multi-view stereo (MVS) $\rightarrow$ Poisson surface reconstruction \cite{luhmann2014photogrammetry,lowe2004sift,furukawa2010pmvs,kazhdan2006poisson}. It then describes depth fusion with a truncated signed distance function (TSDF) as used in real-time RGB-D and LiDAR systems \cite{newcombe2011kinectfusion}, which underpins ARKit-style scene meshes that RoomPlan builds upon. A comparison between image-only object reconstruction and depth-assisted room understanding clarifies why Object Capture and RoomPlan behave differently. Finally, the chapter maps these ideas to Apple's Object Capture and RoomPlan APIs and documents the 'CapturedRoom' data model.


\section{Motivation}
\label{sec:motivation}

Users redesign interiors more confidently when they can see alternatives in context. Traditional representations (2D plans, catalog photos, or material swatches) omit depth, occlusion, and the relation between existing objects and new ones. Augmented reality can composite virtual furniture into a live video stream, but that mode ties the experience to being physically present and holding the device. If the room is first captured as a structured 3D model, the user can reopen and edit that model later: the scan becomes a persistent spatial substrate for experimentation.

Research datasets such as ScanNet \cite{dai2017scannet} demonstrated that indoor applications benefit from joint geometry and semantics (which surfaces are walls, which objects are chairs or tables). RoomPlan brings a consumer subset of that idea to supported iPhones and iPads: it outputs not only geometry but labelled entities with rough dimensions. This thesis uses that output as the basis for furniture replacement, substituting USDZ models for detected instances, which goes beyond merely placing a free-floating object on a manually identified floor plane.

\section{Photogrammetry and the Classical Reconstruction Pipeline}
\label{sec:photogrammetry}

\subsection{Definition and triangulation}

\textbf{Photogrammetry} is the discipline of extracting metric information from photographs \cite{luhmann2014photogrammetry}. In 3D reconstruction, the elementary geometric operation is 
triangulation: corresponding points in two or more calibrated images determine a 3D point by intersecting rays. Stable triangulation requires baseline between cameras and sufficient overlap so that the same surface points appear in multiple views.
\begin{figure}[htbp]
\centering
% ch2_photogrammetry_concept.tex
% fig:ch2_photo_concept — photogrammetry triangulation concept

\begin{tikzpicture}[
  cam/.style={rectangle, rounded corners=3pt, draw=black, fill=gray!12,
              minimum width=2.6cm, minimum height=1.0cm,
              font=\small\bfseries, align=center, line width=0.7pt},
  pt/.style={circle, draw=black, fill=black, minimum size=6pt, inner sep=0pt},
  obs/.style={circle, fill=gray!60, draw=black, minimum size=5pt, inner sep=0pt},
  ray/.style={dashed, draw=gray!70, line width=0.8pt},
  ba/.style={rectangle, rounded corners=3pt, draw=black, fill=white,
             minimum width=5.4cm, minimum height=1.5cm,
             font=\small, align=center, line width=0.7pt},
  arrow/.style={-{Stealth[length=5pt]}, thick, dashed, draw=gray!80}
]

%% ── 3D world point ──
\node[pt] (X) at (5.5, 7.2) {};
\node[font=\small\bfseries, above=5pt of X] {$\mathbf{X}$\ (3D point)};

%% ── Cameras ──
\node[cam] (C1) at (0.5,  3.5) {Camera 1\\$\mathbf{P}_1$};
\node[cam] (C2) at (5.5,  3.0) {Camera 2\\$\mathbf{P}_2$};
\node[cam] (C3) at (10.5, 3.5) {Camera 3\\$\mathbf{P}_3$};

%% ── Rays from cameras to X ──
% C1 top-right corner → X
\draw[ray] (C1.north east) -- (X);
% C2 top → X
\draw[ray] (C2.north)      -- (X);
% C3 top-left corner → X
\draw[ray] (C3.north west) -- (X);

%% ── Observed image points placed ON each ray ──
% Each point is at ~35% of the way from camera to X.
% C1=(1.8,4.0) → X=(5.5,7.2): at t=0.35 → (1.8+0.35*3.7, 4.0+0.35*3.2) = (3.1, 5.1)
\node[obs, label={[font=\footnotesize]left:$\mathbf{x}_1$}]
  (x1) at (3.1, 5.1) {};

% C2=(5.5,4.0) → X=(5.5,7.2): at t=0.35 → (5.5, 4.0+0.35*3.2) = (5.5, 5.1)
\node[obs, label={[font=\footnotesize]right:$\mathbf{x}_2$}]
  (x2) at (5.5, 5.1) {};

% C3=(9.2,4.0) → X=(5.5,7.2): at t=0.35 → (9.2-0.35*3.7, 4.0+0.35*3.2) = (7.9, 5.1)
\node[obs, label={[font=\footnotesize]right:$\mathbf{x}_3$}]
  (x3) at (7.9, 5.1) {};

%% ── Triangulation label ──
\node[font=\footnotesize\itshape, text=gray!80] at (5.5, 6.3)
  {Triangulation: intersect rays $\;\rightarrow\;$ recover $\mathbf{X}$};

%% ── Bundle adjustment box ──
\node[ba] (ba) at (5.5, 0.9) {%
  \textbf{Bundle Adjustment}\\[4pt]
  $\displaystyle
  \min_{\mathbf{P}_j,\,\mathbf{X}_i}
  \sum_{i,j}
  \bigl\|\mathbf{x}_{ij} - \pi(\mathbf{P}_j,\mathbf{X}_i)\bigr\|^2$
};

\node[font=\footnotesize\itshape, text=gray, below=0.1cm of ba]
  {jointly refines all camera poses and 3D points};

%% ── Arrows from BA to each camera — all identical dashed style ──
\draw[arrow] (ba.west)  -- (C1.south);
\draw[arrow] (ba.north) -- (C2.south);
\draw[arrow] (ba.east)  -- (C3.south);

\end{tikzpicture}

\caption{Triangulation and Bundle Adjustment: corresponding image points from calibrated cameras determine a 3D point~$\mathbf{X}$; Bundle Adjustment jointly minimises reprojection error over all camera poses~$\mathbf{P}_j$ and scene points~$\mathbf{X}_i$.}
\label{fig:ch2_photo_concept}
\end{figure}

\subsection{From features to textured mesh}

A widely used multi-view pipeline proceeds in stages \cite{luhmann2014photogrammetry,lowe2004sift,furukawa2010pmvs,kazhdan2006poisson}:

\begin{enumerate}
    \item \textbf{Feature detection and matching}: Distinctive keypoints (e.g.\ SIFT \cite{lowe2004sift}) are extracted and matched across images.

    \item \textbf{Structure from Motion}: Camera intrinsics and extrinsics are estimated together with a \emph{sparse} 3D point cloud. \textbf{Bundle adjustment} refines parameters by minimising reprojection error across all observations. Symbolically, one minimises $\sum_{i,j} \|\mathbf{x}_{ij} - \pi(\mathbf{P}_j, \mathbf{X}_i)\|^2$ over projection matrices $\mathbf{P}_j$ and points $\mathbf{X}_i$, where $\mathbf{x}_{ij}$ are measured image points and $\pi$ is the projection function.

    \item \textbf{Multi-view stereo}: Dense correspondence and depth estimation across views yield a dense point cloud or per-view depth maps \cite{furukawa2010pmvs}.

    \item \textbf{Surface reconstruction}: The dense samples are converted to a watertight or oriented surface, commonly using \textbf{Poisson surface reconstruction} \cite{kazhdan2006poisson}, which solves for an indicator function whose gradient aligns with surface normals, then extracts a mesh (e.g.\ Marching Cubes).

    \item \textbf{Texture mapping}: The mesh is parameterised; colours are blended from source images onto the surface.
\end{enumerate}
\begin{figure}[htbp]
\centering
% fig2_2_sfm_mvs.tex
% fig:ch2_photo_pipeline — SfM–MVS vertical flowchart (fully sequential)

\begin{tikzpicture}[
  node distance=0.7cm,
  boxtop/.style={rectangle, rounded corners=6pt,
                 draw=red!40, fill=red!6, line width=0.8pt,
                 minimum width=5.5cm, minimum height=1.0cm, align=center},
  boxmid/.style={rectangle, rounded corners=6pt,
                 draw=gray!60, fill=gray!8, line width=0.8pt,
                 minimum width=5.5cm, minimum height=1.0cm, align=center},
  boxout/.style={rectangle, rounded corners=6pt,
                 draw=red!30, fill=gray!15, line width=0.8pt,
                 minimum width=5.5cm, minimum height=0.9cm, align=center},
  arr/.style={-{Stealth[length=5pt]}, thick, draw=gray!60},
  stg/.style={font=\scriptsize\itshape, text=gray!70}
]

\node[boxtop] (imgs) {
  \textbf{\small Multi-view RGB images}\\[1pt]
  {\scriptsize\color{gray!70} overlapping photographs, ARKit poses}
};

\node[boxmid, below=0.7cm of imgs] (feat) {
  \textbf{\small Feature detection and matching}\\[1pt]
  {\scriptsize\color{gray!70} SIFT keypoints extracted and matched across views}
};

\node[boxmid, below=0.7cm of feat] (sfm) {
  \textbf{\small Structure from Motion (SfM)}\\[1pt]
  {\scriptsize\color{gray!70} sparse point cloud, bundle adjustment}
};

\node[boxmid, below=0.7cm of sfm] (mvs) {
  \textbf{\small Multi-view Stereo (MVS)}\\[1pt]
  {\scriptsize\color{gray!70} dense depth maps per view}
};

\node[boxmid, below=0.7cm of mvs] (surf) {
  \textbf{\small Surface Reconstruction}\\[1pt]
  {\scriptsize\color{gray!70} Poisson meshing and texturing}
};

\node[boxout, below=0.7cm of surf] (usdz) {
  \textbf{\small Textured mesh output (USDZ)}
};

\draw[arr] (imgs.south) -- node[stg, right=3pt]{stage 1} (feat.north);
\draw[arr] (feat.south) -- node[stg, right=3pt]{stage 2} (sfm.north);
\draw[arr] (sfm.south)  -- node[stg, right=3pt]{stage 3} (mvs.north);
\draw[arr] (mvs.south)  -- node[stg, right=3pt]{stage 4} (surf.north);
\draw[arr] (surf.south) -- node[stg, right=3pt]{stage 5} (usdz.north);

\end{tikzpicture}

\caption{SfM--MVS reconstruction pipeline: from multi-view image input
through feature detection, Structure from Motion (SfM), dense Multi-View
Stereo (MVS), to surface reconstruction and textured mesh output. Pipeline
stages follow the standard formulation in~\cite{luhmann2014photogrammetry,lowe2004sift,furukawa2010pmvs,kazhdan2006poisson}.}
\label{fig:ch2_photo_pipeline}
\end{figure}

This end-to-end path is what Section~\ref{sec:on_device_pipeline} relates to \textbf{Object Capture}.

\section{Depth Fusion and TSDF-Based Scene Meshing}
\label{sec:depth_fusion}

Room-scale tracking on mobile devices fuses depth (LiDAR or stereo) with visual-inertial odometry. A standard volumetric representation is the TSDF: for each voxel in a bounded volume, one stores a (truncated) signed distance to the nearest surface along with a weight; each new depth frame updates voxels along viewing rays \cite{newcombe2011kinectfusion}. Over time, noise averages down and a coherent surface emerges. The mesh shown in AR debugging views is typically an isosurface of that field.

This reconstruction is driven by measured depth: textureless walls still contribute if the depth sensor returns reliable values. By contrast, MVS in photogrammetry relies on photometric consistency and can fail on untextured or specular surfaces. Semantic interpretation (which planar patch is a ``wall'' and which blob is a ``sofa'') is not solved by the TSDF alone; RoomPlan adds learning-based detection and classification on top of ARKit scene understanding.

\section{Image-Only versus Depth-Assisted Reconstruction}
\label{sec:depth_vs_photo}

\textbf{Object Capture} is primarily an image-driven pipeline (with poses from ARKit): geometry is inferred from many RGB views through stages analogous to SfM, MVS, and Poisson (or similar) meshing, followed by texturing. \textbf{Room Plan} is depth- and semantics-first: it targets a structured, parametric room description ('CapturedRoom') rather than a single closed mesh of the whole interior. Table~\ref{tab:depth_vs_photo} contrasts the two paradigms at the level relevant to this thesis.

\begin{table}[htbp]
\centering
\caption{Comparison of image-centric object reconstruction (Object Capture) and depth-assisted parametric room capture (Room Plan). The Room Plan branch builds on real-time depth integration in the TSDF family \cite{newcombe2011kinectfusion}, augmented by learned structure and semantics.}
\label{tab:depth_vs_photo}
\begin{tabular}{|p{0.28\textwidth}|p{0.30\textwidth}|p{0.30\textwidth}|}
\hline
\textbf{Aspect} & \textbf{Object Capture} & \textbf{Room Plan} \\
\hline
Primary sensing & Dense RGB coverage & LiDAR depth + RGB + IMU \\
\hline
Dominant geometry cue & Photo-consistency (MVS) & TSDF fusion + plane fitting \\
\hline
Typical output & Textured triangle mesh (USDZ) & Parametric walls, openings, objects \\
\hline
Semantics & Optional masking & Built-in categories (e.g.\ chair, table) \\
\hline
\end{tabular}
\end{table}

\section{On-Device Capture Pipeline}
\label{sec:on_device_pipeline}

Figure~\ref{fig:capture_branches} compares two parallel pipelines.  
The left column is the photogrammetric/Object Capture path:
multi-view RGB images and poses $\rightarrow$ feature matching $\rightarrow$ SfM
(sparse structure) $\rightarrow$ MVS (dense geometry) $\rightarrow$ surface reconstruction
and texturing $\rightarrow$ USDZ output.
The right column is the RoomPlan path:
LiDAR/RGB/IMU input $\rightarrow$ VIO tracking $\rightarrow$ TSDF/depth fusion and mesh anchors
$\rightarrow$ plane/structure inference $\rightarrow$ semantic object/opening detection
$\rightarrow$ parametric fitting $\rightarrow$ 'CapturedRoom' output.
Apple’s internal implementations are proprietary; therefore, the figure
is a conceptual mapping of stages to public APIs.

% \begin{table}[htbp]
% \centering
% \caption{Conceptual mapping from classical stages to Apple's on-device APIs.}
% \label{tab:pipeline_to_api}
% \begin{tabular}{|l|p{0.35\textwidth}|p{0.35\textwidth}|}
% \hline
% \textbf{Stage} & \textbf{Object Capture} & \textbf{Room Plan} \\
% \hline
% Camera trajectory & ARKit VIO; session stores images & ARKit world tracking during scan \\
% \hline
% Sparse / dense geometry & Photogrammetry engine (SfM/MVS class) & TSDF-class fusion + parametric fitting \\
% \hline
% Surface model & Textured mesh & Structural elements + object boxes \\
% \hline
% Semantics & Background masking (when enabled) & Wall / door / window / object classes \\
% \hline
% Serialised output & USDZ file & 'CapturedRoom' (Codable) \\
% \hline
% \end{tabular}
% \end{table}

\begin{figure}[H]
\centering
% fig2_3_dual_pipeline.tex
% fig:capture_branches — two vertical pipelines side by side

\begin{tikzpicture}[
  node distance=0.6cm,
  %% Shared stage style
  stage/.style={rectangle, rounded corners=5pt, draw=gray!60, fill=gray!8,
                line width=0.8pt, minimum width=4.8cm, minimum height=0.95cm,
                font=\small, align=center},
  %% Input boxes
  input/.style={stage, draw=red!40, fill=red!6},
  %% Output / terminal boxes
  terminal/.style={stage, draw=gray!80, fill=gray!22, font=\small\bfseries},
  %% Header labels
  header/.style={font=\small\bfseries, align=center},
  %% Arrows
  arr/.style={-{Stealth[length=5pt]}, thick, draw=gray!60},
  %% Italic annotations
  ann/.style={font=\footnotesize\itshape, text=gray!70, align=center}
]

%%────────────────────────────────────────
%% LEFT COLUMN — Object Capture
%%────────────────────────────────────────
\node[header] (hA) at (0, 0) {Object Capture};

\node[input,    below=0.5cm of hA]  (imgs) {Multi-view RGB Images};
\node[stage,    below=0.6cm of imgs](feat) {Feature Detection \& Matching};
\node[stage,    below=0.6cm of feat](sfm)  {Structure from Motion (SfM)};
\node[stage,    below=0.6cm of sfm] (mvs)  {Multi-View Stereo (MVS)};
\node[stage,    below=0.6cm of mvs] (surf) {Surface Reconstruction (Poisson)};
\node[terminal, below=0.6cm of surf](usdz) {USDZ Mesh};

\draw[arr] (imgs) -- (feat);
\draw[arr] (feat) -- (sfm);
\draw[arr] (sfm)  -- (mvs);
\draw[arr] (mvs)  -- (surf);
\draw[arr] (surf) -- (usdz);

%%────────────────────────────────────────
%% RIGHT COLUMN — RoomPlan
%%────────────────────────────────────────
\node[header] (hB) at (6.2, 0) {RoomPlan};

\node[input,    below=0.5cm of hB]  (sens) {LiDAR + RGB + IMU};
\node[stage,    below=0.6cm of sens](vio)  {Visual-Inertial Odometry (VIO)};
\node[stage,    below=0.6cm of vio] (tsdf) {TSDF Depth Fusion};
\node[stage,    below=0.6cm of tsdf](plan) {Plane Fitting \& Semantics};
\node[stage,    below=0.6cm of plan](param){Parametric Fitting};
\node[terminal, below=0.6cm of param](room){\texttt{CapturedRoom} JSON};

\draw[arr] (sens) -- (vio);
\draw[arr] (vio)  -- (tsdf);
\draw[arr] (tsdf) -- (plan);
\draw[arr] (plan) -- (param);
\draw[arr] (param)-- (room);

%%────────────────────────────────────────
%% Shared output — RealityKit Editor
%%────────────────────────────────────────
\node[terminal, minimum width=4.0cm,
      below=1.4cm of $(usdz)!0.5!(room)$] (ed) {RealityKit Editor};

\draw[arr] (usdz.south) |- (ed.west);
\draw[arr] (room.south) |- (ed.east);

%%── Role annotations ──
\node[ann, below=0.12cm of usdz] {furniture assets (Asset A)};
\node[ann, below=0.12cm of room] {room substrate};

%%── Vertical divider between columns ──
\draw[dashed, gray!40, line width=0.7pt] (3.1, 0.3) -- (3.1, -9.8);

\end{tikzpicture}

% \thesisfig{1\textwidth}{figures/fig2_3_dual_pipeline.pdf}
\caption{Parallel capture pipelines: Object Capture follows an 
image-based reconstruction chain (left column); RoomPlan combines 
LiDAR depth fusion with parametric and semantic inference to produce 
a 'CapturedRoom' (right column).}
\label{fig:capture_branches}
\end{figure}

\subsection{Object Capture}

ObjectCaptureSession guides the user around an object so that views overlap; PhotogrammetrySession reconstructs a mesh and materials, outputting USDZ. That format is well supported in RealityKit and matches the furniture assets used in the editor.

\subsection{Room Plan and CapturedRoom}

Room Plan exposes types such as walls, doors, windows, and CapturedRoom.Object instances with 'category', 'dimensions', and 'transform'. Apple continues to extend taxonomy and behaviour \cite{apple2024roomplan2}. Because 'CapturedRoom' conforms to 'Codable', applications can persist scans as JSON or similar, as done in the implementation (Chapter~\ref{ch:implementation}) under a per-room file in the Documents hierarchy.

\subsection{Feature-to-API mapping}

Table~\ref{tab:feature_api} ties user-visible capabilities to frameworks.

\begin{table}[h]
\centering
\caption{Major application features and the APIs or artefacts that implement them.}
\label{tab:feature_api}
\begin{tabular}{|l|l|p{0.42\textwidth}|}
\hline
\textbf{Feature} & \textbf{API / artefact} & \textbf{Role} \\
\hline
Room scanning & RoomPlan & Produce 'CapturedRoom' \\
\hline
Room persistence & 'Codable' + file I/O & Reload scans after app restart \\
\hline
3D editor & RealityKit (non-AR) & Visualise and edit room \\
\hline
Object scanning & Object Capture & Author USDZ furniture \\
\hline
Splat viewing & Metal + '.ply' & Compare with meshes \\
\hline
\end{tabular}
\end{table}

\section{Limitations}
\label{sec:limitations_capture}

Photogrammetry remains sensitive to glass, mirrors, and uniform surfaces. Room Plan can miss small objects, merge adjacent pieces, or mislabel categories; coarse bounding boxes are not CAD precision. These issues motivate manual touch-up in the editor and, for research-oriented quality, optional export to Blender for refinement; parametric RoomPlan geometry is exported as glTF boxes and markers, with USDZ furniture applied via post-import script, including 3DGS training pipelines discussed in Chapter~\ref{ch:theory2}.  

